\begin{table}[h]
\small
%\scriptsize
\centering
\caption{Example instructions of Bioinfomatics and Compuational Chemistry tasks (Section \ref{sec:data_collection}).}
\begin{tabular}
{
@{\hspace{0pt}}lc@{\hspace{0pt}}
}
\toprule
\textbf{Disciplines} & \textbf{Task Instructions}\\
\midrule

\multirow{19.5}{*}{Bioinformatics} & Train a cell counting model on the BBBC002 datasets containing\\
& Drosophila KC167 cells. Save the test set predictions as a single \\
& column ``count'' to ``pred\_results/cell-count\_pred.csv''. \\
\cmidrule{2-2}

& Train a drug-target interaction model using the DAVIS dataset to \\ 
& determine the binding affinity between several drugs and targets. \\
& Then use the trained model to predict the binding affinities between \\ 
& antiviral drugs and COVID-19 target. Rank the antiviral drugs based \\ 
& on their predicted affinities and save the ordered list of drugs to \\
& ``pred\_results/davis\_dti\_repurposing.txt'', with one SMILES per line. \\
\cmidrule{2-2}

& Plot the Tanimoto similarities of the fingerprint between the frames. \\
& Specifically, the interaction fingerprints between a selected ligand \\
& and protein for the first 10 trajectory frames. Save the png file \\ 
& into pred\_results/ligand\_similarity\_pred.png. \\
\cmidrule{2-2}

& Train a VAE model on the given data and perform a 1-vs-all \\
& differential expression test for each cell type. Extract top markers \\
& for each cell type using the results. Visualize them as a dotplot with \\
& the cell types organized using a dendrogram. Save the figure to \\
& pred\_results/hca\_cell\_type\_de.png. \\

\midrule

\multirow{20.5}{*}{Computational Chemistry} & Train a multitask model on the Clintox dataset to predict a drug's \\
& toxicity and FDA approval status. Save the test set predictions, \\
& including the SMILES representation of drugs and the probability \\
& of positive labels, to ``pred\_results/clintox\_test\_pred.csv''. \\
\cmidrule{2-2}

& Generate features for the given diffusion data based on material \\
& composition and use the SHAP feature selection approach to select \\
& 20 features. Save the selected features as a CSV file \\
& ``mat\_diffusion\_features.csv'' to the folder ``pred\_results/''. \\
\cmidrule{2-2}

& Filter the compounds in ``hits.csv'' and save the SMILES represen- \\
& tations of the left ones. Compounds to be kept should have no PAINS \\
& or Brenk filter substructures and have a maximum tanimoto similarity \\
& of less than 0.5 to any of the active compounds in ``train.csv''. Save \\
& the SMILES of left compounds to ``pred\_results/compound\_filter\_results \\
& .txt'', with each one in a line. \\
\cmidrule{2-2}

& Train a graph convolutional network on the given dataset to predict \\
& the aquatic toxicity of compounds. Use the resulting model to compute \\
& and visualize the atomic contributions to molecular activity of the \\
& given test example compound. Save the figure as \\
& ``pred\_results/aquatic\_toxicity\_qsar\_vis.png''.\\

\bottomrule
\end{tabular}
\vspace{-11pt}
\label{tab:biochem_task_inst}
\end{table}

\newpage

\begin{table}[h]
\small
%\scriptsize
\centering
\caption{Example instructions of Geographical Information Science and Psychology \& Cognitive Neuroscience tasks (Section \ref{sec:data_collection}).}
\begin{tabular}
{
@{\hspace{0pt}}lc@{\hspace{0pt}}
}
\toprule
\textbf{Disciplines} & \textbf{Task Instructions}\\
\midrule

\multirow{20.5}{*}{Geo Information Science} & Analyze and visualize Elk movements in the given dataset. Esti- \\ 
& mate home ranges and assess habitat preferences using spatial  \\
& analysis techniques. Identify the spatial clusters of Elk move- \\ 
& ments. Document the findings with maps and visualizations.  \\
& Save the figure as ``pred\_results/Elk\_Analysis.png''. \\
\cmidrule{2-2}

& Analyze the impact of land subsidence on flooding based on \\
& future elevation data of the study area. Identify flood-prone \\ 
& areas and estimate potential building damage to support urban \\
& planning and mitigation strategies. Save the results to \\
& ``pred\_results/flooding\_analysis.png''. \\
\cmidrule{2-2}

& Calculate the deforestation area percentage in the Brazilian \\ 
& state of Rondônia within the buffer zone of 5.5km around \\
& road layers. Save the percentage result in a CSV file named \\
& ``pred\_results/deforestation\_rate.csv'' with a column title \\
& percentage\_deforestation. \\
\cmidrule{2-2}

& Load North America climate data in NetCDF file and extract \\
& temperature data along the time series, then perform a \\
& quadratic polynomial fit analysis on the temperature data, \\
& and output the fitting results by year in \\
& `pred\_results/polynomial\_fit\_pred.csv'. \\

\midrule

\multirow{25.5}{*}{Psy \& Cognitive Neuroscience} & Process and visualize the given ECG data by perform R \\
& peak detection and outlier correction. Plot an overview of \\ 
& the data and save the final figure as \\
& ``pred\_results/ecg\_processing\_vis1\_pred\_result.png''.\\
\cmidrule{2-2}

& Analyze the inertial measurement unit (IMU) data collected \\
& during sleep and compute sleep endpoints. Load the given data \\
& and compute the following sleep endpoints: time of falling asleep, \\ 
& time of awakening, and total duration spent sleeping. The three \\
& values should be saved in a JSON file ``pred\_results/imu\_ \\
& pred.json'', and the keys for them are "sleep\_onset", \\
& "wake\_onset", and "total\_sleep\_duration", respectively. \\
\cmidrule{2-2}

& Analyze cognitive theories using pattern similarity. Process CSV \\
& files containing model predictions for various syllogistic \\
& reasoning tasks. Calculate similarity scores between these models \\
& and pre-computed high-conscientiousness and high-openness patterns. \\
& The results will contain similarity scores for each cognitive model \\
& with respect to the personality trait patterns. Save the results to \\
& ``pred\_results/CogSci\_pattern\_high\_sim\_data\_pred.csv''. \\
\cmidrule{2-2}

& Train a linear model to learn the mapping of neural represen- \\
& tations in EEG signals from one subject (Sub 01) to another \\
& (Sub 03) based on the preprocessed EEG data from Sub 01 and \\
& Sub 03. Then use the test set of Subject 1 (Sub 01) to gene- \\
& rate EEG signal of Subject 3 (Sub 03). Save the generated EEG \\
& signal of Subject 3 to ``pred\_results/linear\_sub01tosub03\_pred.npy''.\\

\bottomrule
\end{tabular}
\vspace{-11pt}
\label{tab:geopsy_task_inst}
\end{table}